% !TeX TS-program = xelatex
\documentclass[aspectratio=169, 12pt, xcolor=table]{beamer}
\usefonttheme{professionalfonts}
\usefonttheme{serif}
\usepackage[T1]{fontenc}
\usepackage{fontspec-xetex}

\usepackage{tikz}

\usepackage{booktabs}
\usepackage{ifthen}
\usepackage{listings}
\usepackage{subcaption}

\usetikzlibrary{shapes.geometric, arrows}

\setmainfont{Lato}

%\PassOptionsToPackage[more=table]{xcolor}

% Local configuration
\renewcommand{\figurename}{}
\DeclareCaptionFormat{custom}
{%
	\tiny #3
}
\captionsetup{format=custom}

% Title stuff
\title{Programming}
\subtitle{Advanced Python Fun: Part I}
\date{Week 10}
\author{Vassilis Markos, Mediterranean College}

\usetheme{streamline}

% Local Commands
\newcommand{\ohref}[1]{\href{#1}{\texttt{#1}}}
\newcommand{\listindex}[2]{{\underset{#1}{\small #2}}}

% Code listings

\definecolor{codegreen}{rgb}{0,0.6,0}
\definecolor{codegray}{rgb}{0.5,0.5,0.5}
\definecolor{codepurple}{rgb}{0.58,0,0.82}
\definecolor{backcolour}{rgb}{0.95,0.95,0.92}

\lstdefinestyle{mystyle}{
	backgroundcolor=\color{backcolour},   
	commentstyle=\color{codegreen},
	keywordstyle=\color{magenta},
	numberstyle=\tiny\color{codegray},
	stringstyle=\color{codepurple},
	basicstyle=\ttfamily\footnotesize,
	breakatwhitespace=false,         
	breaklines=true,                 
	captionpos=b,                    
	keepspaces=true,                 
	numbers=left,                    
	numbersep=5pt,                  
	showspaces=false,                
	showstringspaces=false,
	showtabs=false,                  
	tabsize=2
}

\lstset{style=mystyle}

% Tikz style

\tikzstyle{startstop} = [ellipse, rounded corners, minimum width=2cm, minimum height=0.8cm, text centered, draw=black, fill=none]
\tikzstyle{io} = [trapezium, trapezium left angle=70, trapezium right angle=110, minimum width=2cm, minimum height=0.8cm, text centered, draw=black, fill=none]
\tikzstyle{process} = [rectangle, minimum width=2cm, minimum height=0.8cm, text centered, draw=black, fill=none]
\tikzstyle{decision} = [diamond, minimum width=2cm, minimum height=0.8cm, text centered, draw=black, fill=none]
\tikzstyle{arrow} = [thick,->,>=stealth]

\tikzstyle{simple node} = [rectangle, draw=black, minimum width = 0.8cm, minimum height = 0.6cm]

% makeatletter stuff

\makeatletter
\newcommand{\arabicthree}[1]{\expandafter\@arabicthree\csname c@#1\endcsname}
\newcommand{\@arabicthree}[1]{\ifnum #1<100 0\fi\ifnum #1<10 0\fi\number#1}
\makeatother

\newcounter{exno}
\setcounter{exno}{0}

\newcommand{\exno}{\stepcounter{exno}In--class Exercise \#\arabicthree{exno}}

\begin{document}

	\begin{frame}
		\titlepage
	\end{frame}

	\begin{frame}{Contents}
		\tableofcontents
	\end{frame}

	\section{\texttt{args} And \texttt{kwargs}}\label{sec:args-and-kwargs}
	
	\sectionframe
	
	\begin{frame}{Challenge}
		Define a function \texttt{foo()} that behaves as shown below (DO NOT inspect the contents of foo.py).
		\lstinputlisting[language=Python]{../source/boom_001.py}
	\end{frame}

	\begin{frame}{\texttt{foo}}
		Let us have a look at \texttt{foo()}:
		\lstinputlisting[language=Python]{../source/foo.py}
	\end{frame}

	\begin{frame}{\texttt{*args}}
		\begin{itemize}
			\item \texttt{*args} is a nice and compact way python allows us to pass positional arguments to a function.
			\item Roughly, you can thinkg of \texttt{*args} arguments as a tuple (immutable list) of values that is passed as a single argument into a list.
			\item Essentially, this means that we can provide a function with an arbitrary number of arguments, without having to define multiple versions of it!
			\item Order matters, in the sense that the contents of \texttt{*args} are spread in the same order they are provided to the function!
		\end{itemize}
	\end{frame}

	\begin{frame}{What Will This Print?}
		\lstinputlisting[language=Python]{../source/boom_002.py}
	\end{frame}
	
	\begin{frame}[fragile]{Be Careful!}
		The above should print something like that:
\begin{lstlisting}
3
2
Traceback (most recent call last):
  File "boom_002.py", line 11, in <module>
    print(bar(1))
TypeError: bar() missing 1 required positional argument: 'b'
\end{lstlisting}\pause
		\begin{itemize}
			\item While \texttt{*args} might be an empty tuple, if we also define other positional arguments, we have to pass them no matter what!\pause
			\item So, in this case, \texttt{bar()} should always be called with \textbf{at least two arguments!}
		\end{itemize}
	\end{frame}
	
	\begin{frame}{Antoher Challenge}
		Define a function \texttt{baz()} that behaves as shown below (DO NOT inspect the contents of baz.py).
		\lstinputlisting[language=Python]{../source/boom_003.py}
	\end{frame}

	\begin{frame}{\texttt{baz()}}
		\lstinputlisting[language=Python]{../source/baz.py}
	\end{frame}

	\begin{frame}{Keyword Arguments}
		\begin{itemize}
			\item Much like list arguments, keyword arguments are arguments provided usign keyrwords, in the same form as a Python dictionary.
			\item The same restrictions to spread positional arguments apply, i.e.:
			\begin{itemize}
				\item We can provide an arbitrary number of keyword arguments.
				\item We can iterate them in the same way we do with dictionaries.
				\item Order does not matter, since we are passing arguments by keyword.
			\end{itemize}
			\item That way, one can pass named arguments to a function, using the same names as in the function's definition.
		\end{itemize}
	\end{frame}

	\begin{frame}{\texttt{*args}: The Other Way Around}
		We can also use \texttt{*args} to spread list contents as arguments into a list, by using the list spread operator, \texttt{*}.
		\lstinputlisting[language=Python]{../source/boom_004.py}
	\end{frame}

	\begin{frame}{\texttt{**kwargs}: The Other Way Around}
		We can also use \texttt{**kwargs} to spread dict contents as arguments, by using the dictionary spread operator, \texttt{**}.
		\lstinputlisting[language=Python]{../source/boom_005.py}
	\end{frame}

	\begin{frame}{Some Remarks}
		\begin{itemize}
			\item (Reminder) You cannot use a positional argument after a keyword one!
			\item So, you have to collect all postional arguments at the beginning and all keyword arguments at the end of the function parameters list.
			\item This also extends to \texttt{*args} and \texttt{**kwargs}, which means that we cannot expect for positional arguments after using \texttt{**kwargs} --- in the same spirit as above.
			\item One common use of \texttt{**kwargs} is to provide default values to keyword arguments and then allow users to provide custom values whenever they want to.
		\end{itemize}
	\end{frame}
	
	\section{Fun Time!}\label{sec:fun-time}
	
	\sectionframe
	
	\setcounter{exno}{0}
	
	\begin{frame}{\exno: Revision (Arrays)}
		Solve all exercises found in:
		\begin{center}
			\ohref{labs/arrays\_lab/arrays\_01.pdf}
		\end{center}
		Pay attention to the subtle points of each exercise and what you should compute in each time (especially in the two last ones).
	\end{frame}

	\begin{frame}{\exno: Revision (Files)}
		Solve all exercises found in:
		\begin{center}
			\ohref{labs/files\_lab/files\_01.pdf}
		\end{center}
		In all cases, do not hardcode any paths and use appropriate Python libraries to handle proper path finding.
	\end{frame}
	
	\begin{frame}{\exno: Revision (Classes)}
		Solve all exercises found in:
		\begin{center}
			\ohref{labs/classes\_lab/objects\_01.pdf}
		\end{center}
		As with the two previous ones, stay calm and pay attention to details. Also, make your designs as succinct as possible!
	\end{frame}

	\begin{frame}{Any Questions?}
		\begin{minipage}{0.35\textwidth}
			\raggedright
			Do not forget to fill in the questionnaire shown right!

			\vspace{2\topsep}

			\begin{scriptsize}
				\begin{tabular}{l}
					\textbf{Office Hours} \\
					Glyfada: Tuesday, 09:00--10:00, PC Lab \\
					Pellinis: Wednesday, 15:00--17:00, 203
				\end{tabular}				
			\end{scriptsize}
		\end{minipage}\hfill
		\begin{minipage}{0.58\textwidth}
			\vspace{0pt}
			\raggedleft
			\includegraphics[scale=0.4]{../../assets/post_lesson_assessment.png}
			\centering
			\ohref{https://forms.gle/dKSrmE1VRVWqxBGZA}
		\end{minipage}
	\end{frame}
	
\end{document}